\pdfvariable suppressoptionalinfo 611
\providecommand{\CRevisionRound}{2}
\documentclass[11pt]{article}
\usepackage[a4paper,margin=25mm]{geometry}
\usepackage{amsmath,amssymb,amsthm,booktabs,microtype}
\usepackage[T1]{fontenc}
\usepackage{lmodern}
\usepackage[hidelinks]{hyperref}
\newtheorem{theorem}{Theorem}
\newcommand{\Sreg}{\mathcal S_+}
\title{An Exact Event-Time Atlas for an Affine-Impact Bouncing Ball}
\author{HCS Research Program}
\date{28 August 2026\quad (revision \CRevisionRound)}
\begin{document}
\maketitle
\begin{abstract}
We analyse a gravity-driven ball on a half-line with restitution and a constant
impact impulse.  The physical flow is $\dot q=v$, $\dot v=-g$, the guard is
$q=0,v^-<0$, and the reset is $v^+=r(-v^-)+J$.  On the positive outgoing
section the event map is the affine map $P(u)=ru+J$, but each iterate carries
the physical roof $2u/g$.  This gives exact geometric and arithmetic event
sequences, a sharp Zeno boundary, a unique positive forced flight for
$J>0,r<1$, and separate elastic, sticking, and accelerating boundaries.  We
adjoin $(0,0)$ as rest and exclude it from the regular section, so a
zero-duration fixed event is not counted as a physical cycle.  The certificate
checks 96 exact event-time cells and 36 interior impact reconstructions.  Any
fixed-point series is explicitly an event-map object, not a physical-flow
zeta or a target arithmetic construction.
\end{abstract}

\section{Hybrid model and reduction}
For $g>0$, $0\le r\le1$, and $J\ge0$, let
\[
 \dot q=v,\qquad \dot v=-g\quad(q>0),\qquad
 v^+=r(-v^-)+J\quad(q=0,v^-<0).
\]
The point $(q,v)=(0,0)$ is a separate absorbing rest state.  For an interior
state, solving $q_0+v_0t-gt^2/2=0$ gives
\[
 \tau_0=\frac{v_0+\sqrt{v_0^2+2gq_0}}{g},\qquad
 w_0=\sqrt{v_0^2+2gq_0},\qquad u_0=rw_0+J.
\]
Starting just after an impact with positive outgoing speed $u$, symmetry of
constant acceleration gives one physical flight of duration $2u/g$ and returns
the incoming speed $-u$.  Thus the outgoing section map is
\[
 P(u)=ru+J,\qquad \tau(u)=\frac{2u}{g},\qquad u\in\Sreg=(0,\infty).
\]

\begin{theorem}[event-time atlas]\label{thm:impact}
For $r\ne1$, put $u_*=J/(1-r)$.  Then
\[
 u_n=u_*+r^n(u_0-u_*),\qquad
 t_n=\frac2g\left[n u_*+(u_0-u_*)\frac{1-r^n}{1-r}\right].
\]
For $r=1$, $u_n=u_0+nJ$ and
$t_n=\frac2g[n u_0+Jn(n-1)/2]$.  If $J=0$ and $0<r<1$, positive
speeds form a Zeno sequence with accumulation time
$2u_0/[g(1-r)]$.  The edge $r=0,J=0$ has one positive-duration flight at
most and then sticks at rest.  If $J>0$ and $r<1$, $u_*$ is the unique
positive forced cycle, with physical period $2u_*/g$ and event multiplier
$P'(u)=r$; every positive speed converges to it.  At $r=1,J=0$ there is a
continuum of elastic periods $2u/g$, while $r=1,J>0$ is nonperiodic with
quadratic event times.
\end{theorem}
\begin{proof}
The first equation is the quadratic flight root.  The reset gives $P$, and
induction followed by a finite geometric or arithmetic sum gives the displayed
iterates and times.  The geometric sum is finite precisely for
$J=0,0<r<1$; $r=0$ reaches zero after one flight and the rest convention
stops the execution.  Solving $u=P(u)$ and using contraction proves the forced
cycle and its multiplier.  Identity and translation follow directly.
\end{proof}

\ifnum\CRevisionRound>0
\section{Section domains and formal series}
The regular section is $\Sreg=(0,\infty)$ (the literal label
\texttt{S\_+}): it parametrizes positive-duration
flights only.  Consequently, for $J>0,r<1$ the physical event-map fixed-point
series is $\zeta_{\rm phys}(z)=1/(1-z)$, whereas for $J=0,r<1$ it is the
empty series $1$ on $\Sreg$.  If the section is artificially closed to
$[0,\infty)$, the affine boundary point $u=0$ contributes the separate formal
series $\zeta_{\rm aff}(z)=1/(1-z)$ when $J=0,r<1$; that point is rest, not
a physical orbit.  For $r=1,J=0$ the fixed set is a continuum and a
cardinality series is undefined.  For $r=1,J>0$ translation has no fixed
point, so its finite-count series is $1$.  None of these event-map series is a
physical-flow zeta.
\fi

\ifnum\CRevisionRound>1
\section{Exact audit and scope boundary}
Twelve rational controls cover contraction, strict Zeno, the $r=0$ sticking
edge, elastic identity, and translation.  The ledger contains 96 exact roof and
cumulative-time cells and 36 first-impact reconstructions.  An independent
checker passes 368 exact assertions; SymPy passes 11 identities, byte replay is
exact, and 14 repaired/stale hash mutations are rejected.  In particular, an
attack that relabels $r=0,J=0$ as Zeno is caught before release.

The Route-A tuple is
\begin{center}\small\ttfamily
(A0\_FAIL, A1\_WEAK, A2\_FAIL, A3\_FAIL, A4\_FORMAL\_HINT),
\end{center}
with \texttt{overall=ROUTE\_A\_REJECTED} and
\texttt{route\_b\_invocation\_allowed=false}.  The package makes no target
prime, zero, local factor, root number, automorphy, target functional equation,
or Hilbert--Polya operator claim.
\fi

\section*{Source note}
Hybrid modeling context is given by Leine--Nijmeijer~\cite{leine2004} and
Goebel--Sanfelice--Teel~\cite{goebel2012}; no priority claim is made.
\begin{thebibliography}{9}
\bibitem{leine2004} R. I. Leine and H. Nijmeijer, \emph{Dynamics and
Bifurcations of Non-Smooth Mechanical Systems}, Springer (2004),
DOI:\href{https://doi.org/10.1007/978-3-540-44398-8}{10.1007/978-3-540-44398-8}.
\bibitem{goebel2012} R. Goebel, R. G. Sanfelice, and A. R. Teel,
\emph{Hybrid Dynamical Systems: Modeling, Stability, and Robustness},
Princeton University Press (2012),
DOI:\href{https://doi.org/10.1515/9781400842636}{10.1515/9781400842636}.
\end{thebibliography}

\section*{Declarations}
\textbf{Scope.} The exact registered literal is
\texttt{NO\_BAD\_EULER\_OR\_ROOT\_NUMBER}.
\textbf{Data and code.} The theorem, exact ledger, independent checks, and
build recipe are released with HCS-C212. \textbf{Competing interests.} None
declared. \textbf{AI-use disclosure.} Generative tools assisted drafting and
code generation; the internal artifact chain checked formulas, metadata, and
scope boundaries.  This is not external peer review.
\end{document}
